\documentclass[10pt,a4paper]{amsart}
\usepackage{amsmath,amssymb,amsthm,graphicx}
% Retained minimal preamble for the Stacks Project chapter sites-cohomology.tex (Cohomology on Sites).
% Verbatim excerpt of the chapter preamble: Xy-pic setup (native preamble lines 14-18) and the native
% theorem declarations with the mathematical macros (native preamble lines 154-217). Copied mechanically;
% nothing added, removed or redefined. The wrapper class and the amsmath/amssymb/amsthm/graphicx base
% packages are supplied by the assembler, so the class/externaldocument part of the native preamble is omitted.
% For commutative diagrams we use Xy-pic
\usepackage[all]{xy}

% We use 2cell for 2-commutative diagrams.
\xyoption{2cell}
% Theorem environments.
%
\theoremstyle{plain}
\newtheorem{theorem}[subsection]{Theorem}
\newtheorem{proposition}[subsection]{Proposition}
\newtheorem{lemma}[subsection]{Lemma}

\theoremstyle{definition}
\newtheorem{definition}[subsection]{Definition}
\newtheorem{example}[subsection]{Example}
\newtheorem{exercise}[subsection]{Exercise}
\newtheorem{situation}[subsection]{Situation}

\theoremstyle{remark}
\newtheorem{remark}[subsection]{Remark}
\newtheorem{remarks}[subsection]{Remarks}

\numberwithin{equation}{subsection}

% Macros
%
\def\lim{\mathop{\mathrm{lim}}\nolimits}
\def\colim{\mathop{\mathrm{colim}}\nolimits}
\def\Spec{\mathop{\mathrm{Spec}}}
\def\Hom{\mathop{\mathrm{Hom}}\nolimits}
\def\Ext{\mathop{\mathrm{Ext}}\nolimits}
\def\SheafHom{\mathop{\mathcal{H}\!\mathit{om}}\nolimits}
\def\SheafExt{\mathop{\mathcal{E}\!\mathit{xt}}\nolimits}
\def\Sch{\mathit{Sch}}
\def\Mor{\mathop{\mathrm{Mor}}\nolimits}
\def\Ob{\mathop{\mathrm{Ob}}\nolimits}
\def\Sh{\mathop{\mathit{Sh}}\nolimits}
\def\NL{\mathop{N\!L}\nolimits}
\def\CH{\mathop{\mathrm{CH}}\nolimits}
\def\proetale{{pro\text{-}\acute{e}tale}}
\def\etale{{\acute{e}tale}}
\def\QCoh{\mathit{QCoh}}
\def\Ker{\mathop{\mathrm{Ker}}}
\def\Im{\mathop{\mathrm{Im}}}
\def\Coker{\mathop{\mathrm{Coker}}}
\def\Coim{\mathop{\mathrm{Coim}}}

% Boxtimes
%
\DeclareMathSymbol{\boxtimes}{\mathbin}{AMSa}{"02}

%
% Macros for moduli stacks/spaces
%
\def\QCohstack{\mathcal{QC}\!\mathit{oh}}
\def\Cohstack{\mathcal{C}\!\mathit{oh}}
\def\Spacesstack{\mathcal{S}\!\mathit{paces}}
\def\Quotfunctor{\mathrm{Quot}}
\def\Hilbfunctor{\mathrm{Hilb}}
\def\Curvesstack{\mathcal{C}\!\mathit{urves}}
\def\Polarizedstack{\mathcal{P}\!\mathit{olarized}}
\def\Complexesstack{\mathcal{C}\!\mathit{omplexes}}
% \Pic is the operator that assigns to X its picard group, usage \Pic(X)
% \Picardstack_{X/B} denotes the Picard stack of X over B
% \Picardfunctor_{X/B} denotes the Picard functor of X over B
\def\Pic{\mathop{\mathrm{Pic}}\nolimits}
\def\Picardstack{\mathcal{P}\!\mathit{ic}}
\def\Picardfunctor{\mathrm{Pic}}
\def\Deformationcategory{\mathcal{D}\!\mathit{ef}}

% --- appended 2026-10-12 (Chief): the project own definitions of the reference/slogan/history
% display environments, copied mechanically from the shared Stacks preamble lines 8-12
% (\usepackage{verbatim} plus the three \newenvironment lines, which the chapter preamble keeps
% verbatim from the project preamble). Nothing added, removed or redefined: these environments
% carry no mathematical content and the project itself typesets them as comments. Without these
% three lines a counted statement whose own text uses them fails with Environment slogan undefined.
\usepackage{verbatim}
\newenvironment{reference}{\comment}{\endcomment}
\newenvironment{slogan}{\comment}{\endcomment}
\newenvironment{history}{\comment}{\endcomment}
\usepackage{cleveref}
\crefname{originallemma}{original lemma}{original lemmas}
\Crefname{originallemma}{Original lemma}{Original lemmas}
\makeatletter
\newlabel{stacks-lemma-gerbe-abelian-auts}{{0CJY}{source}{}{sourcepointer-1}{}}
\newlabel{stacks-lemma-gerbe-abelian-auts@cref}{{[originallemma][0][]0CJY}{[1][0][]source}{}{}{}}
\makeatother
\crefname{originaldefinition}{original definition}{original definitions}
\Crefname{originaldefinition}{Original definition}{Original definitions}
\makeatletter
\newlabel{stacks-definition-gerbe}{{06NZ}{source}{}{sourcepointer-2}{}}
\newlabel{stacks-definition-gerbe@cref}{{[originaldefinition][0][]06NZ}{[1][0][]source}{}{}{}}
\makeatother
\crefname{originallemma}{original lemma}{original lemmas}
\Crefname{originallemma}{Original lemma}{Original lemmas}
\makeatletter
\newlabel{lemma-torsors-h1}{{03AJ}{source}{}{sourcepointer-3}{}}
\newlabel{lemma-torsors-h1@cref}{{[originallemma][0][]03AJ}{[1][0][]source}{}{}{}}
\makeatother
\crefname{originallemma}{original lemma}{original lemmas}
\Crefname{originallemma}{Original lemma}{Original lemmas}
\makeatletter
\newlabel{lemma-trivial-torsor}{{03AI}{source}{}{sourcepointer-4}{}}
\newlabel{lemma-trivial-torsor@cref}{{[originallemma][0][]03AI}{[1][0][]source}{}{}{}}
\makeatother
\crefname{originallemma}{original lemma}{original lemmas}
\Crefname{originallemma}{Original lemma}{Original lemmas}
\makeatletter
\newlabel{lemma-cech-spectral-sequence}{{03AZ}{source}{}{sourcepointer-5}{}}
\newlabel{lemma-cech-spectral-sequence@cref}{{[originallemma][0][]03AZ}{[1][0][]source}{}{}{}}
\makeatother
\title{An abelian gerbe has an object over any neighbourhood on which the relevant cohomology vanishes}
\author{The Stacks Project Authors}
\date{stacks-project:sites-cohomology.tex:local-snapshot-20261011; GNU FDL 1.2 or later}
\begin{document}
\maketitle
\noindent\textit{Editorial note.} All delivered mathematical text is the author's own, copied from the pinned Stacks Project chapter file sites-cohomology.tex; no statement or proof step was rewritten, reordered or removed.\par\medskip
\noindent\textit{Editorial note.} The proof is the author's complete argument: the local objects, the triviality of the isom torsors, the cohomological obstruction, its vanishing and the correction of the isomorphisms are all written out; the parenthetical that one computation is omitted concerns the verification that the sections form a two-cocycle, a routine cocycle check, and it is preserved verbatim in the counted span.\par\medskip
\noindent\textit{Editorial note.} Every labelled reference occurring inside the counted spans is carried in assembly.external\_references as a non-load-bearing navigation pointer (same-chapter ones at their definition site, cross-chapter ones at the invocation inside the counted span), so no load-bearing mathematics is replaced by a stub. No bibliography citation occurs inside the counted spans.\par\medskip
\noindent\textit{Editorial note.} The checked chapter list for same-name collisions is the set of stacks chapters staged in this repository: chow, cohomology, divisors, etale-cohomology, flat, more-algebra, more-morphisms, morphisms, spaces-more-morphisms, varieties together with this card's chapter sites-cohomology.\par\medskip
\noindent\textit{Editorial note.} The preamble is the chapter's own retained minimal preamble (sources/stacks-sites-cohomology/retained-preamble.tex), a verbatim excerpt of the native preamble: the Xy-pic setup and the native theorem declarations with the mathematical macros.\par\medskip
\noindent\textit{Allowed standard prerequisites (editorial).}\begin{itemize}
\item Gerbes and their abelian automorphism sheaves. Stacks, Definition stacks-definition-gerbe and Lemma stacks-lemma-gerbe-abelian-auts. Supplies the local objects of the gerbe and the abelian sheaf of automorphisms.
\item Torsors under an abelian sheaf and their triviality. Same chapter, Lemmas lemma-torsors-h1 and lemma-trivial-torsor. Turns the vanishing of the first cohomology on the pieces into isomorphisms between the local objects.
\item The Cech-to-cohomology spectral sequence. Same chapter, Lemma lemma-cech-spectral-sequence. Shows that the comparison map from Cech cohomology in degree two to cohomology is injective when the first cohomology of the pieces vanishes.
\end{itemize}\medskip
\ifdefined\hypertarget\hypertarget{sourcepointer-1}{}\fi
\noindent\textit{Source reference.} Original lemma 0CJY: sources/stacks-sites-cohomology/pinned-original-source.tex, source lines 1669--1669. Navigation pointer invoked at proof line 1669; the counted span contains the pointer only and the cited result is not reproduced inside it, so the pointer is registered as a non-load-bearing navigation stub in the established form for Stacks material. No bibliography citation is involved.\par\medskip
\ifdefined\hypertarget\hypertarget{sourcepointer-2}{}\fi
\noindent\textit{Source reference.} Original definition 06NZ: sources/stacks-sites-cohomology/pinned-original-source.tex, source lines 1694--1694. Navigation pointer invoked at proof line 1694; the counted span contains the pointer only and the cited result is not reproduced inside it, so the pointer is registered as a non-load-bearing navigation stub in the established form for Stacks material. No bibliography citation is involved.\par\medskip
\ifdefined\hypertarget\hypertarget{sourcepointer-3}{}\fi
\noindent\textit{Source reference.} Original lemma 03AJ: sources/stacks-sites-cohomology/pinned-original-source.tex, source lines 291--291. Navigation pointer invoked at proof line 1715; the counted span contains the pointer only and the cited result is not reproduced inside it, so the pointer is registered as a non-load-bearing navigation stub in the established form for Stacks material. No bibliography citation is involved.\par\medskip
\ifdefined\hypertarget\hypertarget{sourcepointer-4}{}\fi
\noindent\textit{Source reference.} Original lemma 03AI: sources/stacks-sites-cohomology/pinned-original-source.tex, source lines 278--278. Navigation pointer invoked at proof line 1717; the counted span contains the pointer only and the cited result is not reproduced inside it, so the pointer is registered as a non-load-bearing navigation stub in the established form for Stacks material. No bibliography citation is involved.\par\medskip
\ifdefined\hypertarget\hypertarget{sourcepointer-5}{}\fi
\noindent\textit{Source reference.} Original lemma 03AZ: sources/stacks-sites-cohomology/pinned-original-source.tex, source lines 1475--1475. Navigation pointer invoked at proof line 1738; the counted span contains the pointer only and the cited result is not reproduced inside it, so the pointer is registered as a non-load-bearing navigation stub in the established form for Stacks material. No bibliography citation is involved.\par\medskip
\begin{lemma}
\label{lemma-existence}
Let $\mathcal{C}$ be a site. Let $p : \mathcal{S} \to \mathcal{C}$
be a gerbe over a site whose automorphism sheaves are abelian.
Let $\mathcal{G}$ be the sheaf of abelian groups constructed
in Stacks, Lemma \ref{stacks-lemma-gerbe-abelian-auts}.
Let $U$ be an object of $\mathcal{C}$ such that
\begin{enumerate}
\item there exists a cofinal system of coverings $\{U_i \to U\}$
of $U$ in $\mathcal{C}$ such that $H^1(U_i, \mathcal{G}) = 0$ and
$H^1(U_i \times_U U_j, \mathcal{G}) = 0$
for all $i, j$, and
\item $H^2(U, \mathcal{G}) = 0$.
\end{enumerate}
Then there exists an object of $\mathcal{S}$ lying over $U$.
\end{lemma}

\begin{proof}
By Stacks, Definition \ref{stacks-definition-gerbe}
there exists a covering $\mathcal{U} = \{U_i \to U\}$
and $x_i$ in $\mathcal{S}$ lying over $U_i$.
Write $U_{ij} = U_i \times_U U_j$. By (1) after refining the covering
we may assume that $H^1(U_i, \mathcal{G}) = 0$ and
$H^1(U_{ij}, \mathcal{G}) = 0$.
Consider the sheaf
$$
\mathcal{F}_{ij} =
\mathit{Isom}(x_i|_{U_{ij}}, x_j|_{U_{ij}})
$$
on $\mathcal{C}/U_{ij}$. Since
$\mathcal{G}|_{U_{ij}} = \mathit{Aut}(x_i|_{U_{ij}})$ we see that there
is an action
$$
\mathcal{G}|_{U_{ij}} \times \mathcal{F}_{ij} \to \mathcal{F}_{ij}
$$
by precomposition. It is clear that $\mathcal{F}_{ij}$ is a
pseudo $\mathcal{G}|_{U_{ij}}$-torsor and in fact a torsor because
any two objects of a gerbe are locally isomorphic.
By our choice of the covering and by
Lemma \ref{lemma-torsors-h1}
these torsors are trivial (and hence have global sections by
Lemma \ref{lemma-trivial-torsor}).
In other words, we can choose isomorphisms
$$
\varphi_{ij} : x_i|_{U_{ij}} \longrightarrow x_j|_{U_{ij}}
$$
To find an object $x$ over $U$ we are going to massage our choice
of these $\varphi_{ij}$ to get a descent datum (which is necessarily
effective as $p : \mathcal{S} \to \mathcal{C}$ is a stack).
Namely, the obstruction to being a descent datum is that the cocycle
condition may not hold. Namely, set $U_{ijk} = U_i \times_U U_j \times_U U_k$.
Then we can consider
$$
g_{ijk} = \varphi_{ik}^{-1}|_{U_{ijk}} \circ \varphi_{jk}|_{U_{ijk}} \circ
\varphi_{ij}|_{U_{ijk}}
$$
which is an automorphism of $x_i$ over $U_{ijk}$. Thus we may and do
consider $g_{ijk}$ as a section of $\mathcal{G}$ over $U_{ijk}$.
A computation (omitted) shows that $(g_{i_0i_1i_2})$ is a $2$-cocycle
in the {\v C}ech complex ${\check C}^\bullet(\mathcal{U}, \mathcal{G})$
of $\mathcal{G}$ with respect to the covering $\mathcal{U}$.
By the spectral sequence of
Lemma \ref{lemma-cech-spectral-sequence}
and since $H^1(U_i, \mathcal{G}) = 0$ for all $i$
we see that ${\check H}^2(\mathcal{U}, \mathcal{G}) \to H^2(U, \mathcal{G})$
is injective. Hence $(g_{i_0i_1i_2})$ is a coboundary
by our assumption that $H^2(U, \mathcal{G}) = 0$.
Thus we can find sections $g_{ij} \in \mathcal{G}(U_{ij})$ such that
$g_{ik}^{-1}|_{U_{ijk}} g_{jk}|_{U_{ijk}}  g_{ij}|_{U_{ijk}} = g_{ijk}$
for all $i, j, k$.
After replacing $\varphi_{ij}$ by $\varphi_{ij}g_{ij}^{-1}$
we see that $\varphi_{ij}$ gives a descent datum on the objects
$x_i$ over $U_i$ and the proof is complete.
\end{proof}

\end{document}
